\documentclass[10pt,a4paper]{amsart}
\usepackage[margin=19mm]{geometry}
\usepackage{amsmath,amssymb,mathtools}
\usepackage[scr=rsfs,cal=euler]{mathalfa}
\usepackage{xfrac}
\usepackage[unicode,hidelinks,bookmarks=false]{hyperref}
\newcommand{\ie}{i.e.\ }
\newcommand{\cf}{cf.\ }
\newcommand{\resp}{respectively\ }
\newcommand{\Subsection}{\subsection}
\providecommand{\nobreakdash}{\nobreak\hbox{-}}
\setlength{\emergencystretch}{2em}
\multlinegap0pt
\usepackage{bm}
\usepackage{bbm}
\usepackage{dsfont}
\usepackage{xspace}
\usepackage{cancel}
\usepackage{mathtools}
\usepackage{amsthm}
\makeatletter
\@ifundefined{theorem}{\newtheorem{theorem}{Theorem}}{}
\@ifundefined{definition}{\newtheorem{definition}[theorem]{Definition}}{}
\@ifundefined{example}{\newtheorem{example}[theorem]{Example}}{}
\@ifundefined{lemma}{\newtheorem{lemma}[theorem]{Lemma}}{}
\makeatother
\makeatletter
\newcommand{\CHs}{\mathrm{CH}^\star}
\newcommand{\Hilb}{\mathrm{Hilb}}
\newcommand{\cE}{\mathcal{E}}
\expandafter\ifx\csname customconj\endcsname\relax
\newenvironment{customconj}[1]
  {\renewcommand\theinnercustomconj{#1}\innercustomconj}
\fi
\expandafter\ifx\csname customcor\endcsname\relax
\newenvironment{customcor}[1]
  {\renewcommand\theinnercustomcor{#1}\innercustomcor}
\fi
\expandafter\ifx\csname customthm\endcsname\relax
\newenvironment{customthm}[1]
  {\renewcommand\theinnercustomthm{#1}\innercustomthm}
\fi
\makeatother
\title[Logarithmic Hilbert schemes of curves as weighted blow-ups a]{Logarithmic Hilbert schemes of curves as weighted blow-ups and their integral Chow rings}
\author{Veronica Arena, Terry Dekun Song}
\date{Source: arXiv:2605.17922v1 (CC BY 4.0)}
\begin{document}
\maketitle

\noindent\emph{Source and licence.} Logarithmic Hilbert schemes of curves as weighted blow-ups and their integral Chow rings. Version arXiv:2605.17922v1 (CC BY 4.0), distributed under the Creative Commons Attribution 4.0 International licence, as recorded in the arXiv licence metadata for this exact version and shown on the abstract page https://arxiv.org/abs/2605.17922v1. Full rights notice: licenses/arxiv-2605.17922-CC-BY-4.0.txt.\par\smallskip

\noindent\textit{Editorial scope.} Counted unit: the Example whose source label is \texttt{None} in the article (source0.tex lines 2214--2216, source label \texttt{None}), delivered together with the article's own complete original proof of it (source0.tex lines 2217--2256, one \texttt{proof} environment of 40 lines). No mathematics is written, repaired or re-derived: statement and proof are the article's own text line for line. Required context retained uncounted: def: doubleindexes (lines 2118--2124); lem:rings (lines 2259--2261). Every definition, display and label of those lines is retained unchanged. Omitted from this package and not redistributed: 1--2117, 2125--2213, 2257--2258, 2262--2462. Nothing inside a retained excerpt is omitted or altered: every retained body is delivered exactly as the article's lines are written, so no omission receipt of any kind accompanies this package. Citations: the retained excerpts carry no citation command at all, so no bibliography entry of the article is transcribed and no reference is redistributed. Reference adaptation: none was needed and none was made; every reference inside the delivered unit resolves inside it, so no unresolved reference or citation is produced. Printed slips kept literal: no printed word, symbol, hypothesis or number of the counted statement or of its proof was altered. Layout and class: the article's own class is \texttt{amsart} with packages including subfiles, euscript, tabu, geometry, xcolor, verbatim, graphicx, amssymb, mathrsfs, amsthm, amsmath, stmaryrd, tikz, tikz-cd; the wrapper is the lane's pinned \texttt{amsart} editorial wrapper at 10pt on A4, which restates the theorem-like environments the retained text needs and adds the article's own macro definitions that the retained text needs (\texttt{\textbackslash CHs}, \texttt{\textbackslash Hilb}, \texttt{\textbackslash cE}), with extra wrapper packages (bm, bbm, dsfont, xspace, cancel, mathtools). The delivered numbering is the unit's own. Assets: no retained span contains a figure environment, a graphics-inclusion command, a PDF-page inclusion, a table or a TikZ picture; no image file is copied into this package, the package declares no asset dependency and no image file is redistributed with it. Licence: version arXiv:2605.17922v1 of this article is distributed under the Creative Commons Attribution 4.0 International licence, as recorded in the arXiv licence metadata for this exact version and shown on the abstract page https://arxiv.org/abs/2605.17922v1. A full rights notice is delivered at licenses/arxiv-2605.17922-CC-BY-4.0.txt. Characters: every retained excerpt is pure ASCII, so the delivered engine, XeLaTeX with its default Latin Modern text font, reports no missing character.
\begin{definition}\label{def: doubleindexes}
     Let $\cE_{n,i}$ and $\epsilon_{n,i}=[\cE_{n,i}]\in \mathrm{CH}^1(\Hilb^n(C|p)_{\leq i-1})$ denote the (class of) exceptional divisor introduced along the weighted blow-up $\Hilb^n(C|p)_{\leq i-1}\to \Hilb^n(C|p)_{\leq i}.$ They specify the divisor of subschemes which has length-$i$ supported on the last bubble. 
    
    In general, for $k \geq i,$ let $\cE_{n,k}$ and $\epsilon_{n,k}\in \mathrm{CH}^1(\Hilb^n(C|p)_{\leq i})$ be the proper transforms of $\cE_{n,k}$ and $\epsilon_{n,k}$ on $\Hilb^n(C|p)_{\leq k},$ which agree with the pullback along the iterated weighted blow-ups $\Hilb^n(C|p)_{\leq i}\to \Hilb^n(C|p)_{\leq k}.$

    Let $\mathcal{Z}_{n,i}\subset \Hilb^{n}(C|p)_{\leq i}$ denote the closed substack where exactly $i$ points lie on the endpoint.
\end{definition}

\begin{lemma}\label{lem:rings}
  Let $\tilde{\phi}: R\to S$ be a surjection of rings, and let $I\subset S$ be an ideal such that $\langle b_1,\dots b_n\rangle = I.$ Pick any $a_i\in R$ such that $\tilde{\phi}(a_i) = b_i,$ then $\tilde{\phi}^{-1}(I) = \langle \ker \tilde{\phi}, a_1,\dots,a_n\rangle.$
\end{lemma}

\begin{example}
  When $i = n-1,$ the indices in the list are $j = n.$ Since $n-j = 0,$ all the $Q_{n-j,k}=0,$ so we get $$\mathrm{CH}^\star(\mathrm{Sym}^n(C))[\epsilon_{n,n}]/(Q_{n,n}(\epsilon_{n,n}),\ker (s_{0n}^*)\epsilon_{n,n}).$$
\end{example}

\begin{proof}

For the purpose of the proof, let us rewrite the presentation with the double indexes as in Definition \ref{def: doubleindexes}, so that we can more easily apply induction. With the refined notation, the desired formula reads 

$$\mathrm{CH}^\star(\Hilb^n(C|p)_{\leq i}) = \frac{\mathrm{CH}^\star(\mathrm{Sym}^n(C))[\epsilon_{n,n},\dots, \epsilon_{n,i+1}]}{\left\{\begin{array}{c}
  \text{for each } j = n, \dots, i+1: Q_{n,j}(\epsilon_{n,n},\dots,\epsilon_{n,j+1}, \epsilon_{n,j}); \ \ker(s_{(n-j)n}^*) \cdot \epsilon_{n,j};\\ \text{for each } k = n-j, \dots, 1: Q_{n-j,k}(\epsilon_{n,n},\dots, \epsilon_{n,j+k+1}, \epsilon_{n,j+k})\cdot \epsilon_{n,n-j},  
  \end{array}\right\}}.$$

  We use lexicographic ordering on $(n, n-i)$ and run induction on this totally ordered set. When $n-i = 0,$ the indexing set of $j$ runs from $n, n+1$ which we treat as the empty set, hence the inductive hypothesis on the row $\{(n,0)\}$ holds.

  We present $\Hilb^n(C|p)_{\leq i},$ as blow-up of $\Hilb^n(C|p)_{\leq (i+1)}$ which, by inductive hypothesis, has a presentation $$\CHs(\Hilb^n(C|p)_{\leq (i+1)}) = \frac{\mathrm{CH}^\star(\mathrm{Sym}^n(C))[\epsilon_{n,n},\dots, \epsilon_{n,i+2}]}{\left\{\begin{array}{c}
  \text{for each } j = n, \dots, i+2: Q_{n,j}(\epsilon_{n,n},\dots,\epsilon_{n,j+1}, \epsilon_{n,j}); \  \ker(s_{(n-j)n}^*) \cdot \epsilon_{n,j};\\\text{for each } k = n-j, \dots, 1: Q_{n-j,k}(\epsilon_{n,n},\dots, \epsilon_{n,j+k+1}, \epsilon_{n,j+k})\cdot \epsilon_{n,n-j}, 
  \end{array}\right\}}$$ 

  Applying Keel's formula, we get \[\mathrm{CH}^\star(\Hilb^n(C|p)_{\leq i}) = \frac{\mathrm{CH}^\star(\Hilb^n(C|p)_{\leq i+1})}{(Q_{n,i+1}(\epsilon_{n,i+1}); \ \ker \iota^*_{i+1,n} \cdot \epsilon_{n,i+1})},\] where $\iota_{i+1,n}: \mathcal{Z}_{n,i+1}\to \Hilb^{n}(C|p)_{\leq i+1}$ is the closed embedding. To perform the inductive step, we determine $\ker \iota^*_{i+1,n}.$
  
  In the following, let $m:=n-(i+1).$ By inductive hypothesis, the substack $\mathcal{Z}_{n,i+1}\cong \Hilb^{m}(C|p)_{\leq 0}$ has presentation as $$\CHs(\Hilb^m(C|p)_{\le 0})=\frac{\mathrm{CH}^\star(\mathrm{Sym}^{m}(C))[\epsilon_{m,m},\dots, \epsilon_{m,1}]}{\left\{\begin{array}{c}
   \text{for each } j = m, \dots, 1: Q_{m,j}(\epsilon_{m,m},\dots,\epsilon_{m,j+1}, \epsilon_{m,j}); \  \ker(s_{(m-j),m}^*) \cdot \epsilon_{m,j};\\\text{for each } k = m-j, \dots, 1: Q_{m-j,k}(\epsilon_{m,m},\dots, \epsilon_{m,j+k+1}, \epsilon_{m,j+k})\epsilon_{m,m-j}
  \end{array}\right\}},$$ and the map $$\iota_{m,n}^*: \mathrm{CH}^\star(\Hilb^n(C|p)_{\leq (i+1)})\to \mathrm{CH}^\star(\Hilb^{m}(C|p)_{\leq 0})$$ extends $$s_{m,n}^*: \mathrm{CH}^\star(\mathrm{Sym}^n(C))\to \mathrm{CH}^\star(\mathrm{Sym}^{m}(C))$$ by sending $\epsilon_{n,j}\mapsto \epsilon_{m, j-(i+1)}.$ By Lemma \ref{lem:rings} below, the kernel of the map is generated by $\ker s_{m,n}^*$ and lifts of the relations in the presentation of $\mathrm{CH}^\star(\Hilb^{m}(C|p)_{\leq 0})$ given above.
  
  We observe that for each $j= m, \dots, 1,$ the restriction map $\iota_{m,n}^*$ sends $$Q_{m, j}(\epsilon_{n, j+(i+1)})(\epsilon_{n,n},\dots, \epsilon_{n, j+i+2})\mapsto Q_{m,j}(\epsilon_{m,j})(\epsilon_{m,m},\dots,\epsilon_{m,j+1}),$$ {$$\ker(s_{(\tilde{n}-j)n}^*)\epsilon_{n,j+(i+1)}\mapsto \ker(s_{(\tilde{n}-j)\tilde{n}}^*)\epsilon_{\tilde{n},j}.$$} 
  Given such a $j,$ for each $k = m-j,\dots, 1,$ $\iota_{m,n}^*$ sends $$Q_{m-j, k}(\epsilon_{n, j+k+(i+1)})(\epsilon_{n,n},\dots, \epsilon_{n, j+k+i+2})\epsilon_{n,n-j}\mapsto Q_{m-j,k}(\epsilon_{m,j+k})(\epsilon_{m,m},\dots, \epsilon_{m,j+k+1})\epsilon_{m,m-j}.$$ Therefore, $\ker \iota_{i+1,n}^*$ is generated by
  
$$\left\{\begin{array}{c}
    \ker s_{m,n}^*;  \text{ for each } j= m, \dots, 1:  Q_{m, j}(\epsilon_{n,n},\dots, \epsilon_{n, j+i+2}, \epsilon_{n, j+(i+1)});  \ker(s_{(m-j)n}^*)\cdot \epsilon_{n,j+(i+1);}\\
    \text{for each }k = m-j,\dots, 1:  Q_{m-j, k}(\epsilon_{n,n},\dots, \epsilon_{n, j+k+i+2}, \epsilon_{n, j+k+(i+1)}) \cdot \epsilon_{n,n-j}
  \end{array} \right\}.$$
  
  The new relations in $\mathrm{CH}^\star(\Hilb^n(C|p)_{\leq i})$ introduced by the weighted blow-up are hence generated by 
  $$\left\{\begin{array}{c}
    Q_{n,i+1}(\epsilon_{n,n},\dots, \epsilon_{n,i+2}, \epsilon_{n,i+1}); \ \ker s_{m,n}^* \cdot \epsilon_{n,i+1};
    \\
     \text{for each }j=m, \dots, 1:  Q_{m, j}(\epsilon_{n,n},\dots, \epsilon_{n, j+i+2}, \epsilon_{n, j+(i+1)})\cdot\epsilon_{n,i+1}; \  \ker(s_{(m-j)n}^*)\epsilon_{n,j+(i+1)}\cdot \epsilon_{n,i+1};\\
    \text{for each }k = m-j,\dots, 1:  Q_{m-j, k}(\epsilon_{n,n},\dots, \epsilon_{n, j+k+i+2}, \epsilon_{n, j+k+(i+1)})\cdot \epsilon_{n,n-j} \cdot \epsilon_{n,i+1}
  \end{array} \right\}$$
  
  We notice that among the above terms,  $Q_{m-j, k}(\epsilon_{n,n},\dots, \epsilon_{n, j+k+i+2}, \epsilon_{n, j+k+(i+1)}) \cdot \epsilon_{n,n-j} \cdot \epsilon_{n,i+1}$ lies in the ideal generated by $Q_{m-j, k}(\epsilon_{n,n},\dots, \epsilon_{n, j+k+i+2}, \epsilon_{n, j+k+(i+1)}) \cdot \epsilon_{n,n-j},$ which is already part of relations in $\mathrm{CH}^\star(\Hilb^n(C|p)_{\leq i+1})$ by inductive hypothesis.
  
  Similarly, $\ker s^*_{(n-j),n}\cdot \epsilon_{n,j+(i+1)}$ is a relation in $\mathrm{CH}^\star(\Hilb^n(C|p)_{\leq i+1}).$ The rest of the relations are precisely the relations corresponding to $j = i+1$ (and $k = m,\dots, 1$) as desired.
\end{proof}

\end{document}
